\section{Introducción y objetivos}

\subsection*{Qué se hizo}

Este documento presenta el caso de estudio 2 de Computación de Alto Rendimiento (HPC). Parte del programa secuencial de multiplicación de matrices del caso de estudio 1 y, sobre él, implementa tres cosas que el enunciado pide:

\begin{enumerate}
    \item \textbf{Concurrencia con OpenMP}, con las directivas elegidas y justificadas a partir de los ejemplos de la especificación oficial.
    \item \textbf{Perfilado de la aplicación} en sus dos ejes mínimos: CPU y memoria.
    \item \textbf{Benchmark en dos máquinas distintas}, con hardware realmente diferente (Intel vs AMD, distinta jerarquía de caché, distinto sistema operativo).
\end{enumerate}

El caso de estudio 1 ya había comparado cuatro estrategias: pthreads, \texttt{fork()} con memoria compartida, optimización de compilación (\texttt{-O3} y variantes) y un kernel transpuesto. Aquí se completa el cuadro con OpenMP como quinta estrategia y, sobre todo, se responde con \textbf{medición} una pregunta que antes era intuición: \textit{¿por qué la transposición de $B$ hace que el kernel sea tan más rápido?}

\subsection*{Objetivos}

\begin{itemize}
    \item Implementar la versión concurrente con OpenMP y justificar técnicamente cada directiva y cláusula utilizada.
    \item Caracterizar la aplicación: dónde se va el tiempo de CPU y qué presión genera sobre la memoria.
    \item Medir el rendimiento en dos máquinas Linux distintas y comparar los resultados.
    \item Integrar las tres mediciones (perfilado, speedup y benchmark) en una explicación única del comportamiento del programa.
\end{itemize}

\subsection*{Organización}

Las secciones 2 y 3 presentan los conceptos y el programa base. La sección 4 describe la implementación con OpenMP. La sección 5 expone el diseño experimental. La sección 6 muestra los resultados, la 7 los analiza y la 8 cierra con las conclusiones.
